% Report fragment. Requires booktabs, longtable, array, hyperref, amsmath, amssymb.
\section{Comparison of the complete current PA design corpus}
\label{sec:corpus-comparison}

The comparison is pinned to ProveIt commit
\nolinkurl{793c68737d94bffa4e32a75bce6cc8058d6d22f0}.
It covers all 93 filed research packets and the four additional archives in
\texttt{docs/incoming}: 97 deliveries in total.
The filed-packet descriptions are attributed claims from their authors;
the repository explicitly says their placement did not rerun or verify them.
The archive ledger and graph catalog are independent kinds of evidence.
\footnote{Pinned sources: \href{https://github.com/VladimirReshetnikov/ProveIt/blob/793c68737d94bffa4e32a75bce6cc8058d6d22f0/Computability/BusyBeaver/Research/PA/docs/README.md}{filed packet ledger}; \href{https://github.com/VladimirReshetnikov/ProveIt/blob/793c68737d94bffa4e32a75bce6cc8058d6d22f0/Computability/BusyBeaver/Research/PA/docs/MACHINES.md}{machine census and its equivalence convention}; \href{https://github.com/VladimirReshetnikov/ProveIt/blob/793c68737d94bffa4e32a75bce6cc8058d6d22f0/Computability/BusyBeaver/Research/PA/README.md}{PA project proof status}.}

The supplied \texttt{packet\_ledger97.csv} and JSON counterpart preserve every
delivery, its path, wrapper, reported result and pinned source URL.
The 83 named principal or supplementary graph classes are separately listed
in \texttt{named\_positive\_designs.csv}.  The full structural frontier,
including rejected and auxiliary tables, is in
\texttt{frontier194\_all521.csv}; every one of its 521 representatives is
supplied for exact replay.  These counts concern the input corpus, before
the machines introduced in this report.

\subsection{Counts, equivalence conventions and evidence boundaries}

The official filed catalog records 15,317 original table files, 1,084 distinct
file payloads and 747 graph classes across all state counts.  Of its 194-control
classes, 139 have a spanning root and one is a nonspanning diagnostic.
All 139 spanning historical classes are represented again among the current
incoming comparisons.  Parsing the four current archives adds 381 further
194-control classes, giving 520 spanning classes and the one diagnostic.
Most additions are auxiliary search inputs rather than proposed theorems.

\begin{center}
\begin{tabular}{lr}
\toprule
Role of a current 194-control graph class & Count\\
\midrule
Reported principal construction & 76\\
Reported supplementary construction & 7\\
Historical research auxiliary & 47\\
Current incoming research auxiliary & 372\\
Rejected candidate & 14\\
Research seed & 4\\
Nonspanning diagnostic & 1\\
\midrule
Total & 521\\
\bottomrule
\end{tabular}
\end{center}

The table is deliberately a graph census.  It does not assert 521, or even
83, independently verified consistency equivalences.  A named construction
may be backed by an ordinary PA argument, a relative simulation, a written
\(\mathrm{iPRA}\) argument, or only a narrower artifact claim; the per-packet
ledger retains that distinction.  In particular, the early AX11H delivery
explicitly declined an unconditional \(\mathrm{iPRA}\) theorem, whereas the
second delivery with the same archive name supplies the revised written
calibration used by later packets.  The two deliveries must not be
conflated, and that written calibration is not a completed Lean endpoint.
\footnote{Compare the pinned \href{https://github.com/VladimirReshetnikov/ProveIt/blob/793c68737d94bffa4e32a75bce6cc8058d6d22f0/Computability/BusyBeaver/Research/PA/docs/pa194_ax11h_research_2026-11-14/README.md}{first AX11H delivery} with the \href{https://github.com/VladimirReshetnikov/ProveIt/blob/793c68737d94bffa4e32a75bce6cc8058d6d22f0/Computability/BusyBeaver/Research/PA/docs/pa194_ax11h_research_2026-11-14_2/README.md}{revised AX11H delivery}.}

The direct current-archive scan finds 1,278 \texttt{.tm} paths and 694
byte-distinct payloads.  Recursing into nested archives gives 1,757 table
occurrences and 880 distinct table payloads across six distinct ZIP payloads
(nine ZIP references), of which 677 distinct payloads have 194 controls.
The accompanying JSON audit examines 847 occurrences and 641 distinct JSON
payloads.  Its explicit table recognizers accept array, dictionary and
multiline forms, including Boolean direction encodings; no extra
194-control class appears in the nested archives or recognized JSON tables.
These are separate occurrence, byte and graph counts and must not be added
together.  The scan does not claim to execute every possible abstract
generator or recover tables specified only in prose.

Graph classes disregard state labels and, for the corpus classification,
the designated start.  For novelty of the new machines we use a stronger
exact comparison: arbitrary state renaming and spanning root, reflection,
global symbol exchange, and ignored write/move actions on terminal edges.
The comparison uses exact rooted transition tuples, with hashes only to
bind supplied files.  Global symbol exchange is a structural test here;
it does not itself preserve the convention of a zero blank tape.
The nonspanning diagnostic cannot match a new spanning table.

\subsection{The 31 mechanism families}

The following table collects all 31 identified families.  Its first 27 rows
organize the filed and repeatedly reprinted designs; its last four rows
cover the four current incoming packets.  Names separated by an equality
sign are aliases or an exact graph duplication, not extra constructions.
The primary pairs \((67/7)\) and \((68/8)\) denote AX readers and simultaneous
folds: \(134+67-7=194=134+68-8\).

There are three substantive levels.  Constructor and logical-head changes
alter the numerical calculus and require proof transport.  MP, cache and
AX11 changes alter the command interface and require a proof compiler valid
for arbitrary carried stores.  Register placement, decoder, selector and
shortcut changes preserve an established source or command semantics and
require a complete physical simulation, including blank initialization.
The last column names the relevant proof obligation; it is not a claim
that every archived report discharged it in \(\mathrm{iPRA}\).

\begingroup
\footnotesize
\setlength{\tabcolsep}{3pt}
\renewcommand{\arraystretch}{1.12}
\begin{longtable}{@{}>{\raggedright\arraybackslash}p{.16\linewidth}>{\raggedright\arraybackslash}p{.24\linewidth}>{\raggedright\arraybackslash}p{.30\linewidth}>{\raggedright\arraybackslash}p{.24\linewidth}@{}}
\toprule
Family & Named designs & Recurring behavior & Proof obligation\\
\midrule
\endfirsthead
\toprule
Family & Named designs & Recurring behavior & Proof obligation\\
\midrule
\endhead
\midrule
\multicolumn{4}{r}{Continued on next page}\\
\endfoot
\bottomrule
\endlastfoot
1. Original source / allocation & Primary PA194 (67/7); PA194 (68/8) & Same original command protocol and separate diagonal calculus. & Complete source simulation and PA calibration. \\
2. Constructor coefficient & M4R; M5; M6 & Coefficients four, five and six change numerical formula coding. & Hereditary proof recoding, including the zero alias. \\
3. Cached MP & PA194C; PA194CS\allowbreak{} & Successful MP stores the consequent in cache and clears visible top; Shift tests cache. & Adjusted all-store command compiler. \\
4. Visible no-push MP & PA194N\allowbreak{}; PA194N\allowbreak{}5; PA194N\allowbreak{}6 & Successful MP retains (b,w,b,0), omitting dummy-zero push; original sum guard. & Reserved-zero handling and arbitrary carried stores. \\
5. Cache-only guard & PA194H\allowbreak{}; PA194H\allowbreak{}5; PA194H\allowbreak{}6; PA194H\allowbreak{}S = filed PA194N\allowbreak{}S & Visible no-push MP with Shift test h=1; HS uses symmetric comparison. & Final contradiction detection must use Pop. \\
6. Symmetric sum guard & PA194N\allowbreak{}S (sum guard); PA194SUMS & Original sum condition t+h=1 with symmetric destructive comparison. & Source tail equivalence; NS and SUMS differ in MP push behavior. \\
7. Q distribution & PA194Q\allowbreak{} & Exchange the curried premises of the universal-distribution head. & Restricted-rho derivation of premise exchange. \\
8. Q with coefficient changes & PA194Q\allowbreak{}5; PA194Q\allowbreak{}6 & Q head plus coefficient five or six; original dummy-push MP. & Both premise exchange and radix recoding. \\
9. Q with no-push MP & PA194Q\allowbreak{}N = PA194N\allowbreak{}Q; PA194Q\allowbreak{}N5; PA194Q\allowbreak{}N6 & Q head, visible no-push MP, original sum guard. & Composition of Q proof transport and no-push compiler. \\
10. Q with cache-only guard & PA194Q\allowbreak{}H; PA194Q\allowbreak{}H5; PA194Q\allowbreak{}H6; PA194Q\allowbreak{}HS & Q head, no-push MP, cache-only Shift. & Cache-only Pop compiler and Q transport. \\
11. Q with symmetric sum guard & PA194Q\allowbreak{}NS; PA194Q\allowbreak{}S & Q head with the original sum guard; QNS removes MP push, QS retains it. & Exact sum-guard block equivalence. \\
12. Diagonal choice & PA194D; PA194DHS; Gamma(D) + CSZA & Use Gamma(D) instead of Gamma(A); accompanying controller variants. & New diagonal selector; supplement has narrower evidence. \\
13. Known-zero source returns & PA194A; PA194F; PA194J; PA194ZR; PA194R; PA194CS\allowbreak{}Z = PA194X; PA194CZ; PA194H\allowbreak{}A; PA194H\allowbreak{}AZ; CSZA final-register 3; CSZA final-register 5; CSZA final-register 6 & Replace selected jumps or returns by decrements of registers proved zero. & A source invariant and separately verified blank startup are both needed. \\
14. Physical register placement & PA194-AD; PA194F\_AD; PA194R\_AD; PA194TDY; PA194P5; PA194H\allowbreak{}X & Preserve logical commands while exchanging register positions and routers. & Register-frame conjugacy and complete initialization. \\
15. Decoder topology & PA194B; PA1944L & Preserve the original source program. & All-tape or complete source-frame positive simulation. \\
16. AX11 interface & PA194AX11\allowbreak{} & Push the guard and fall through the full MP operand reader: four operands, eleven entries. & All eleven selectors and the fourth operand must be compiled. \\
17. AX11 with cache-only guard & PA194AX11\allowbreak{}H; PA194AX11\allowbreak{}QH & Four-operand AX11 with H guard, with original or Q distribution. & Final Pop compiler; latest incoming supplies expanded full iPRA claim. \\
18. AX11 with changed coding & PA194AX11\allowbreak{}Q6; PA194AX11\allowbreak{}QH5 & AX11 and Q with coefficient six/sum guard or five/cache guard. & Four-field coding, radix transport, appropriate halt detector. \\
19. AX11 with symmetric sum guard & PA194AX11\allowbreak{}S & Exactly the AX11 sum-guard command semantics. & Relative source/physical equivalence. \\
20. Explicit generalization binder & PA194G; PA194AX11\allowbreak{}G & The first AX operand c controls U(c,h), instead of inherited offset d; ten- or eleven-entry AX. & Shorter generalization word with no offset-setting command. \\
21. Explicit binder with H guard & PA194GH; PA194GHS & Explicit binder and ten-entry AX with cache-only Shift. & New all-store compiler; complete physical frame checks. \\
22. AX11QH decoder alternatives & PA194AX11\allowbreak{}QH4L; PA194AX11\allowbreak{}QH4R & Exactly the existing four-operand AX11QH source. & Positive finite-time transfer to the parent. \\
23. AX11 explicit binder and cache guard & PA194AX11\allowbreak{}GH = PA194AX11\allowbreak{}QGH; PA194AX11\allowbreak{}GHS & Four-field eleven-entry AX, direct binder c, visible no-push MP and cache-only Shift; GHS symmetric comparator. & Explicit all-store proof compiler, four-field scheduler and complete physical/source proof; absolute calibration inherited from revised AX11H manuscript. \\
24. Exact physical acceleration & PA194AX11\allowbreak{}GHF; PA194AX11\allowbreak{}GHSF & Preserves GH/GHS commands; replaces two physical paths of length3 by one transition. & Three finite predecessor invariants, 12 local cases; primitive-recursive clock t$\leq$$\sigma$(t)$\leq$3t. \\
25. Loop continuation reassociation & PA194LC\allowbreak{}; PA194LC\allowbreak{}\_A; PA194LC\allowbreak{}\_B; PA194LC\allowbreak{}\_C; PA194LC\allowbreak{}\_AX11H & Preserves resolved source program, changes implementation of recurring constructor loops. & Source identity, literal-frame contracts and positive polynomial clock; revised AX11H absolute calibration for LC\_AX11H. \\
26. Simultaneous history selectors & PA194SEL; single/composite selector auxiliary classes & Merges decoder checkpoints and recycles former labels as selectors reading last written bit. & Full local-tape identities with 1/3-step expansion and primitive-recursive relative iPRA clock. \\
27. Optimized nine-register placement & PA194AX11\allowbreak{}HR & Preserves AX11H source and arithmetic engines. & All 4,778 literal contracts and whole-tape initialization; relative internal physical simulation. \\
28. Independent binder and diagonal with ten-entry AX & PA194GD; PA194GDH; PA194GDHS & Explicit generalization binder c plus diagonal Gamma(D), ordinary distribution; ten-entry AX; sum/H/HS detectors. & Relative primitive-recursive simulation/command proof plus explicitly inherited revised AX11H calibration; no completed Lean endpoint. \\
29. Loop continuation with two accelerations & PA194LC\allowbreak{}AX11GHF; PA194LC\allowbreak{}AX11GHSF & Loop-continuation reassociation plus two guarded shortcuts; separate tagged-three calculus prototype reaches207 states, not 194. & Relative primitive-recursive simulation/command proof plus explicitly inherited revised AX11H calibration; no completed Lean endpoint. \\
30. Loop continuation with third guarded return shortcut & PA194LC\allowbreak{}G; PA194LC\allowbreak{}GS & GH/GHS loop-continuation reassociation and three guarded shortcuts including native190/1. & Relative primitive-recursive simulation/command proof plus explicitly inherited revised AX11H calibration; no completed Lean endpoint. \\
31. Compatible selector and acceleration composition & PA194AX11\allowbreak{}H\_SF; PA194LC\allowbreak{}\_AX11H\_SF & One selector pair127/124 and two guarded accelerations on AX11H/LC\_AX11H. & Relative primitive-recursive simulation/command proof plus explicitly inherited revised AX11H calibration; no completed Lean endpoint. \\
\end{longtable}
\endgroup

The family grouping is a comparison of mechanisms, not a total ordering of
proof strength.  For example, a cached-only Shift guard changes the final
contradiction-extraction word; a symmetric implementation of the same guard
requires a different low-level comparator proof but no new abstract
calculus.  The explicit binder \(c\) and diagonal choice \(\Gamma(d+x)\)
are independent axes.  The new AX11GDH and AX11GDHS constructions combine
those axes with the eleven-entry, four-operand interface; their ordinary
universal-distribution head is kept fixed.

\subsection{The four current incoming additions}

The nine current principal constructions extend the 74 older named
principal/supplementary graph classes to 83.  All nine are reported with
194 working states, and none of their packets reports an accepted smaller
blank-start machine:
\begin{itemize}
\item \href{https://github.com/VladimirReshetnikov/ProveIt/blob/793c68737d94bffa4e32a75bce6cc8058d6d22f0/docs/incoming/pa194_gd_research_2026-11-14.zip}{GD, GDH and GDHS}. Independent binder and offset diagonal with ten-entry AX; sum, cache-only and symmetric cache-only detectors.
\item \href{https://github.com/VladimirReshetnikov/ProveIt/blob/793c68737d94bffa4e32a75bce6cc8058d6d22f0/docs/incoming/pa194_lcfast_tagged_research_2026-11-14.zip}{LCAX11GHF and LCAX11GHSF}. Loop-continuation reassociation combined with two guarded physical shortcuts. The separate tagged-three-calculus prototype has 207 states.
\item \href{https://github.com/VladimirReshetnikov/ProveIt/blob/793c68737d94bffa4e32a75bce6cc8058d6d22f0/docs/incoming/pa194_loop_return_2026-11-14.zip}{LCG and LCGS}. The preceding loop-continuation family with a third guarded return shortcut.
\item \href{https://github.com/VladimirReshetnikov/ProveIt/blob/793c68737d94bffa4e32a75bce6cc8058d6d22f0/docs/incoming/pa194_comparison_and_new_designs_2026-11-15.zip}{AX11H\_SF and LC\_AX11H\_SF}. A compatible history-selector pair combined with the two previously established guarded accelerations.
\end{itemize}
The reports in these four archives compose relative simulations and command
arguments with the specifically identified revised AX11H written
calibration.  Their reuse of that calibration must remain visible; it
cannot be replaced by a graph hash or a finite blank-tape run.

\subsection{Rejected smaller machines and the remaining gap}

The filed census contains 54 historical 193-state graph classes.  They
include rejected candidates and runtime cores and do not constitute 54
accepted blank-start machines.  The current incoming material reuses the
known Runtime193 and AX11 Runtime193 cores.  The strongest supplied startup
search already permits arbitrary values in all nine primary registers;
simply relaxing a zero-proof-store requirement would repeat an existing
search.  That search reports 767,923 overlapping runs and no complete
admissible frame.  Its geometric, path and spare-list requirements are part
of the acceptance condition, so reaching a familiar control state or an
initial-looking prefix is insufficient.

Consequently the input corpus supplies many mechanisms that can be
combined, but no accepted sub-194 blank-start construction.  The new finite
searches described elsewhere in this report probe additional combinations;
their failure to find a smaller machine is not a lower bound.
